
Selecting 7.8\% of training data for GRPO code generation yields a 4.$24\times$ gradient signal advantage over random selection --- yet that advantage is entirely invisible during training when the model never generates a single correct solution. This discrepancy between a confirmed analytical advantage and an empirically null training result is the central finding of this paper.

Reinforcement Learning from Execution Feedback (RLEF) has emerged as an effective approach for improving code generation capabilities in large language models \citep{Gehring2024}. In RLEF, a policy model receives binary execution rewards (pass/fail on test cases) and is updated via Group Relative Policy Optimization (GRPO) \citep{Shao2024}. Binary rewards impose a fundamental inefficiency: when a model either passes every completion or fails every completion within a group, the within-group reward standard deviation equals zero, producing zero gradient. At group size G=4, this zero-gradient pathology affects approximately 69.25\% of training groups \citep{Nie2026} --- the majority of gradient steps are wasted.

The natural remedy is data selection: train only on problems where the model occasionally succeeds. Several recent methods achieve this via online selection --- computing difficulty signals during training and adjusting the data distribution as the model improves [Jiang et al., 2026; Baroian \& Berger, 2026; Sun et al., 2025; Cui et al., 2026]. Online methods work but face a practical limitation: selection overhead is coupled to training cost, and the selection signal cannot be reused across hyperparameter sweeps or model variants.

This paper proposes offline variance-guided selection: profile a frozen model once on all candidate problems, compute per-problem binary execution reward variance v\_i = p\_i($1-p_{i})$, and select problems with the highest variance before training begins. This design separates a cheap profiling phase (approximately 32 seconds for 374 problems via vLLM) from the expensive GRPO training phase, enabling reuse of the selection across runs.

Three gaps in existing work motivate this study: no prior work has (1) characterized the per-problem binary execution reward variance distribution across MBPP for a code LLM systematically, (2) compared offline variance-guided selection to random selection under binary RLEF, or (3) identified the cold-start precondition that determines whether any selection method can operate in this setting.

The key finding, confirmed experimentally, is: offline frozen-model variance profiling produces a statistically robust selection signal, but GRPO gradient concentration can only manifest this advantage when the model is already capable of generating at least some correct solutions. Cold-start --- the model generating zero correct solutions across all training conditions --- renders data selection method irrelevant. When every problem yields k=0 correct completions in every group, $\sigma =\sqrt(k(G-k))/G=0$ regardless of which problems were selected.

This paper reports this negative result alongside the confirmed analytical findings because the precise characterization of where and why the method fails prevents others from repeating the same failure. The analytical validity of the selection signal (confirmed, 4.$24\times$ advantage) is separable from the training-phase failure (cold-start due to parameter mismatch) --- making that separation explicit is itself a contribution.

\textbf{Contributions:}

\begin{enumerate}
\item \textbf{Variance distribution characterization} (empirical, confirmed): The first systematic per-problem characterization of binary execution reward variance across MBPP training problems for DeepSeek-Coder-7B-Instruct under constrained generation. The distribution is severely right-skewed: 91.7\% of problems yield p\_i=0 (all-fail), 7.8\% yield p\_i=0.25, and 0.5\% yield p\_i=1.0.
\end{enumerate}

\begin{enumerate}
\item \textbf{Offline variance profiling} (method, analytically validated): Top-50 selection by frozen-model variance achieves 4.$24\times$ higher mean variance than random-50 selection (Mann-Whitney U p=2.$22\times 10^{-6})$, with the variance-selected subset stochastically dominating random-50. vLLM batch inference completes the 374-problem profiling in approximately 32 seconds on H100 NVL.
\end{enumerate}

\begin{enumerate}
\item \textbf{Cold-start diagnosis} (negative result, high confidence): More than 50,000 GRPO generation attempts across 50--200 training steps --- all producing zero reward. Four candidate root causes are enumerated and ranked; generation parameter mismatch is identified as the most likely explanation. A prescriptive resolution protocol is provided but not yet empirically validated.
\end{enumerate}

